\documentclass[10pt]{article}

\usepackage[utf8]{inputenc}

\usepackage[T1]{fontenc}

\usepackage{amsmath}

\usepackage{amsfonts}

\usepackage{amssymb}

\usepackage[version=4]{mhchem}

\usepackage{stmaryrd}

\usepackage{graphicx}

\usepackage[export]{adjustbox}

\graphicspath{ {./images/} }

\usepackage{mathrsfs}

\usepackage{bbold}



\begin{document}

П. ф. используется и в статистич. физике. Напр., введем $s$-частичные функции распределения $N$-частичной системы:



\[

\begin{gathered}

F_{s}\left(t, x_{1}, \ldots, x_{n}\right)=V^{s} \int d x_{s+1} \ldots d x_{N} w_{N}\left(t, x_{1}, \ldots, x_{n}\right), \\

s=1,2, \ldots,

\end{gathered}

\]



где $V$ - объем, $x_{i}=\left(\boldsymbol{q}_{i}, \boldsymbol{p}_{i}\right)$, а полная функция распределения $w_{N}$ Удовлетворяет Лиувилля уравнению



\section{с Гамильтона функцией}

\[

d w_{N} / d t=\left\{\boldsymbol{H}, w_{N}\right\}

\]



\[

H=T\left(p_{i}\right)+\boldsymbol{U}\left(\left|q_{i}-q_{j}\right|\right)

\]



Тогда всю цепочку Боголюбова уравнений для $F_{S}$ порождает (в термодинамич. пределе $V, N \rightarrow \infty, V / N=v=$ const) уравнение



\[

\begin{gathered}

\frac{\partial F}{\partial t}=d x f(x)\left\{T, \frac{\delta F}{\delta f(x)}\right\}+\int \frac{1}{2} \int d x d y\left(f ( x ) \left(f(y)+v^{-1} f(x)+\right.\right. \\

\left.+v^{-1} f(y)\right)\left\{U, \frac{\delta^{2} F}{\delta f(x) \delta f(y)}\right\},

\end{gathered}

\]



для П.ф.



$F[t ; f]=\int d x_{1} \ldots d x_{N} w_{N}\left(t, x_{1}, \ldots, x_{N}\right) \prod_{1 \leqslant i \leqslant N}\left(1+v f\left(x_{i}\right)\right)$, а сами $F_{S}$ выражаются через него формулами



\[

F\left(t ; x_{1}, \ldots, x_{s}\right)=\prod_{1 \leqslant i \leqslant s}\left(1-\frac{i}{N}\right)^{-1} \frac{\delta^{2} F}{\delta f\left(x_{1}\right) \ldots \delta f\left(x_{s}\right)} .

\]



Лит.: [1] Бе р езин Ф. А., Метод вторичного квантования, 2 изд., М., 1986; [2] В а с и ль е в А. Н., Функциональные методы в квантовой теории поля и статистике, Л., 1976; [3] С ла в но в А. А., Ф а д д е е в Л. Д., Введение в квантовую теорию калибровочных полей, 2 изд., М., 1988; [4] И ц и к с о н К., 3 ю б е р Ж.-Б., Квантовая теория поля, пер. с англ., т. 1-2, М., 1984.



А. М. Малокостов, В. П. Павлов.



ПРОМЕХУТОЧНЫЕ БОЗОНЫ, $W^{ \pm}$- и $Z^{0}-б$ о зо ны короткоживущие частицы со спином 1 и массами $m_{W} \approx 80$ ГэВ и $m_{Z} \approx 90$ ГэВ соответственно, обмен которыми приводит к слабому взаимодействию между кварками и лептонами. При этом большая масса П.б. объясняет контактный характер этого взаимодействия. Обнаружение $W$ и $Z$-бозонов в соударениях на ускорителе ЦЕРН в 1983 К. Руббиа и С. Ван дер Меером (см. [1]) и совпадение их основных характеристик с предсказаниями теории стало одним из наиболее ярких свидетельств в пользу единой теории слабых и электромагнитных взаимодействий. $W$-бозон распадается по каналам $e v_{l}, \mu v_{\mu}, \tau v_{\tau}, q \bar{q}, Z$-бозон по каналам $e^{+} e^{-}, \mu^{+} \mu^{-}, \tau^{+} \tau^{-}, \bar{q} q$. Ширина распадов П.б. составляет примерно 2 ГэВ.



Лит.: [1] Рубб и а К., «Успехи физич. наук», 1985, т. 147, № 2, $\begin{array}{ll}\text { c. } 371-404 . & \text { А. Ю. Иенатьев. } \\ & \end{array}$



nPONATATOP, Функция распространения, причинная функция Грина, в квантовой теории поля - функция, характеризуюшая распространение релятивистского поля (или его кванта) от одного акта взаимодействия до другого. П. является решением классич. волнового уравнения с $\delta$-образной правой частью, удовлетворяюшим специфич. краевым условиям. Простейший П. $D^{c}(x-y)$ скалярного поля $\varphi(x)$ описывает распространение скалярной частицы между точками пространства-времени $x$ и $у$ и может быть представлен в виде 4-мерного интеграла Фурье:



\[

\begin{gathered}

D^{c}(x-y)=(2 \pi)^{-4} \int e^{i k(x-y)} \Delta^{c}(k) d^{4} k, \\

\Delta^{c}(k)=\left(m^{2}-k^{2}-i \varepsilon\right)^{-1}, \quad \varepsilon \rightarrow+0 .

\end{gathered}

\]



Бесконечно малая мнимая добавка $i \varepsilon$, отвечающая упомянутым выше краевым условиям, дает правило обхода полюсов $\Delta^{c}(k)$, так что после выполнения интегрирования П. оказывается представимым в виде:



\[

D^{c}(x-y)=\Theta\left(x^{0}-y^{0}\right) D^{-}(x-y)-\Theta\left(y^{0}-x^{0}\right) D^{+}(x-y) .

\]



Таким образом, при $x^{0}>y^{0}$ он совпадает с отрицательночастотной частью перестановочной функции Паули - Йордана, равной вакуумному среднему $D^{-}(x-y)=i\langle\varphi(x) \varphi(y)\rangle_{0}$, а при $x^{0}<y^{0}-$ положительно-частотной частью, то есть $i\langle\varphi(y) \varphi(x)\rangle_{0}$. Поэтому



\[

D^{c}(x-y)=i\langle T \varphi(x) \varphi(y)\rangle_{0},

\]



где $T$ - символ хронологического произведения; при $x^{0}>y^{0}$ описывает распространение скалярного кванта из $y$ в $x$, а при $x^{0}<y^{0}-$ из $x$ в $y$. Важность П. в квантовой теории поля связана с тем, что он является основным понятием ковариантной теории возмущений и фигурирует в правилах Фейнмана. Центральная роль П. в квантовополевой теории возмушений впервые установлена Д. Ривье (D. Rivier) и Э. Штюкельбергом (Е. Stueckelberg).



Функцию распространения, учитывающую радиационные поправки при движении частицы между точками $x$ и $y$, называют одетым про пагат ором, или двухточеч ной функцией Грина.



См. также Грина функция в квантовой теории поля, Фейнмана диаграмма.



Jum.: [1] Rivier D., Stueckelberg E., «Phys. Rev.», 1948, v. 74, p. 218; [2] Fe у nman R. P., там же, 1949, v. 76, p. 749; [3] е го же, там же, р. 769; [4] Бо голюбов Н. Н., Ш и рков Д. В., Квантовые $\begin{array}{ll}\text { поля, М., 1993. } & \text { Д. В. Ширков. }\end{array}$ ПPOПATATOP уравнения Лиувилля - функция Грина Лиувилля уравнения для вектора распределения. Вектор распределения $\vec{f}(t)=\left\{f\left(x_{1}, \ldots, x_{s} ; t\right)\right\}$, являющийся совокупностью $s$-частичных функций распределения, удовлетворяет обобщенному уравнению Лиувилля



\begin{center}

\includegraphics[max width=\textwidth]{2023_07_08_92e9610c5e4ea89c35d9g-1}

\end{center}



где $\mathscr{L}-$ оператор Лиувилля. Решение этого уравнения выражается через П. $U(t)$ в виде



\[

\vec{f}(t)=U(t) \vec{f}(0) .

\]



При этом П. уравнения Лиувилля можно определить как решение уравнения



\[

\partial_{t} U(t)=\varphi U(t)

\]



с начальным условием $U(0)=\mathbb{1}$, где 1 - единичный оператор.



Лит.: [1] Б а л е с к у Р., Равновесная и неравновесная статистическая механика, пер. с англ., т. 2, М., 1978.



IPOСАЧИВАНИЕ с Л Учай ного по ля - понятие теории перколяции, характеризующее структуру множеств уровня скалярного однородного относительно трансляций cлучайного поля $\{\xi(x)\}$, где $x \in \mathbb{R}^{d}$, либо $x$ является элементом (вершиной, ребром) периодического графа $G$. Определения П. случайного поля в дискретном $(x \in G)$ и непрерывном $\left(x \in \mathbb{R}^{d}\right)$ случаях несколько отличаются друг от друга. Введем для каждого $h \in \mathbb{R}$ множества уровня $A_{h}^{+}=\{x: \xi(x) \geqslant h\}$, $\mathcal{h}^{-}=\{x: \xi(x)<h\}, x \in G$, либо $\mathbb{R}^{d}$. В дискретном случае поле $\{\xi(x) ; x \in G\}$, по определению, п р о с а ч и в а е т с я, если существует уровень $h^{+}$(п р о с а ч и в а н и я) такой, что при $h>h^{+}$множество $A_{h}^{+}$с вероятностью 1 распадается на конечные связные компоненты, а множество $A_{h}^{-}$с ненулевой вероятностью содержит бесконечную связную компоненту. Рассматривается иногда и двойственное понятие, к-рое эквивалентно просачиванию поля $\{-\xi(x)\}$. В случае непрерывного с вероятностью 1 поля $\left\{\xi(x) ; x \in \mathbb{R}^{d}\right\}$ определение просачивания отличается от определения в дискретном случае требованием сушествования случайной величины $\rho(\xi)$ такой, что множество $A_{h}^{+}$распадалось бы с вероятностью 1 на непересекающиеся компакты, расстояние между к-рыми не меныше $\rho(\xi)$, а множество $A_{h}^{-}$содержало бы с



ПРОСАЧИВАНИЕ 463





\end{document}